\sectionguard
\section*{Appendix: Scope and Qualifications}

\textbf{As of 20 August 2026.} Everything below states the boundaries of this
guide at that date. Retrieval mechanics, checksums, query detail, discarded
leads, the complete passage-by-passage reception matrix, and the
notable-and-quotable and interpretive-proposal audits are in
\texttt{propers/verified.md}, \texttt{research/scope.md} and
\texttt{research/source-bindings.toml} beside this document's source; they are
not repeated here.

\editionnote

\textbf{Edition and formulary identity.} The work is the temporal formulary
printed \latin{DOMINICA DECIMA TERTIA post Pentecosten}, \latin{II classis},
on pp.~393--394 of the 1962 \latin{editio typica} of the \latin{Missale
Romanum} (Vatican City: Typis Polyglottis Vaticanis), with marginal numbers
1562--1571. Green follows the general rubrics for Sundays after Pentecost and
is not printed on those pages; it is recorded here as an inference from the
\latin{Rubricae generales} of 1960, not as a reading of the formulary. The
formulary appoints the \latin{Credo} and the Preface of the Most Holy Trinity,
and appoints no Tract, Sequence, second oration, blessing or ritual text, and
no second Collect, Secret or Postcommunion. Its boundaries are the
Postcommunion of the Twelfth Sunday (no.~1561) and the heading of the
Fourteenth Sunday (no.~1572); nothing between them is omitted.

\textbf{The Sunday number is not part of the prayers' identity.} The Collect
and Postcommunion stand together in the Vatican Gelasian among unnumbered
Sunday Masses and in the Gregorian supplement at \latin{Dominica XIIII}; the
Rheinau and S.\ Gallen Gelasians place the Collect's Mass at \latin{Hebd.\
XV}; Sarum has the set at the Fourteenth Sunday after Trinity. This guide
describes the 1962 formulary and its number, and treats the earlier
assignments as history rather than as alternative identities.

\textbf{Text verification.} All ten elements, their headings, references,
wording, accents, punctuation, rank, both rubrics and both formulary
boundaries were visually collated on 19 August 2026 against the CMAA facsimile
of the Vatican typical edition at 300 and 600 dpi, and independently against
page images of the 1962 Benziger \latin{editio iuxta typicam}. The two
editions differ at the points tabulated in the commentary; those are edition
differences, not doubtful readings, and nothing was blended or substituted.
Three words of the Gospel touched by ink loss in the CMAA scan were resolved
letter by letter against the Benziger images and are printed without brackets
because the readings are secure. The appointed text was separately collated
against the Clementine Vulgate, and the twelve places where it diverges are
tabulated in the commentary as verified textual observations. The Internet
Archive OCR was used only to locate the formulary and controls no published
wording. No \latin{Psalterium Romanum} or other older psalter witness was
collated, so no derivation of any chant reading from a particular psalter is
asserted anywhere in this guide.

\textbf{What this guide is not.} It is a hand-missal companion, not a
liturgical book, a critical edition, a homily, an assembly sheet for one civil
date, or a substitute for the sources cited. It does not treat occurrence,
commemorations, particular calendars, or celebration-mode branches for any
year; those belong to the edition-specific assembly workflow and need their own
evidence.

\textbf{Translation and received text.} The project has composed, translated,
paraphrased, modernised and adapted nothing. English for the scriptural propers
is the Douay--Rheims as revised by Bishop Challoner (1749--1752), the English
of the Clementine Vulgate the missal prints, quoted from the Project Gutenberg
e-text of that revision. English for the three orations is the anonymous
rendering in \work{The Roman Missal translated into the English language for
the use of the laity}, first revised edition (Philadelphia: Eugene Cummiskey,
1861), taken from the three lines that book supplies for this formulary and
reproduced with its own nineteenth-century spelling and its abbreviated
conclusion. Both are public domain in the United States; neither is an official
liturgical translation and neither was made for the 1962 books.

\textbf{Eight places where the registered English does not answer the Latin,
and is not made to.} The Introit's \latin{ne derelínquas}, where the Douay
renders the Clementine's \latin{ne obliviscaris}; the Introit's \latin{voces
quæréntium te}, where the Douay has ``the voices of thy enemies''; the Introit's
and Gradual's \latin{exsúrge, Dómine}, where the Douay has ``Arise, O God''; the
Gradual verse's \latin{memor esto oppróbrii servórum tuórum}, whose Douay
wording stands at Ps.~88:51 and is quoted in the appointed-text apparatus only
as apparatus, never as the verse the missal cites; the Alleluia's \latin{a
generatióne et progénie}; the Offertory's \latin{témpora mea}, where the Douay
reads ``My lots are in thy hands''; the Communion, which is an adapted
centonisation and is therefore shown against the Douay verse printed whole with
its reused phrases marked; and the liturgical incipits \latin{Fratres:} and
\latin{In illo témpore: Dum iret Iesus in Ierúsalem}, which are the missal's
own. Two further gaps fall inside the orations rather than Scripture: the 1861
English collapses \latin{ássequi} in the Collect and drops \latin{augméntum}
and the verb \latin{proficiámus} from the Postcommunion, so both arguments are
conducted on the Latin. In every case the gap is stated where it falls and no
rendering of the project's own is supplied.

\textbf{Search bounds.} Reception was searched in Latin, in Greek, and in
English, across the major reasonably accessible patristic, medieval and later
corpora relevant to each appointed passage: Jerome, Augustine, Chrysostom,
Ambrosiaster, Theodoret, Marius Victorinus, Aquinas and the \work{Glossa
ordinaria} on Galatians; Augustine's \work{Quaestiones Evangeliorum} and
\work{Sermo} 176, Bede, Cyril of Alexandria, Ambrose, Gregory the Great, Bruno
of Segni, Theophylact, the pseudo-Augustinian \work{De vera et falsa
poenitentia} and the \work{Catena aurea} on Luke 17; Augustine's
\work{Enarrationes}, Cassiodorus's \work{Expositio Psalmorum}, Jerome's
\work{Commentarioli} and both series of his \work{Tractatus}, Theodoret's
\work{Interpretatio in Psalmos}, the \work{Glossa ordinaria} on the Psalter, and
Bellarmine's \work{Explanatio in Psalmos} on the three appointed psalms, at the
appointed verses; Rabanus Maurus on Wisdom; the Verona, Gelasian and Gregorian
sacramentaries in their public-domain critical editions; the decrees of Trent;
and the liturgical commentary of Guéranger and Schuster. Consequential negative
results are reported in the commentary rather than concealed --- among them
that Ambrose has no exposition of the appointed Gospel beyond one transitional
clause, that Gregory preached no homily on it although his \work{Moralia}
supplies the reading Bede transcribed, that no Chrysostom homily on it was
found at all, that Augustine's \work{Expositio} on Galatians never treats the
four hundred and thirty years, that Marius Victorinus's comment on Galatians
3:11--19 is lost in a lacuna, that the Secret and Postcommunion are not in
Feltoe's Verona collection, that Wisdom 16:20's wording does not occur in
Ambrose's \work{De mysteriis}, that \latin{fac nos amare quod praecipis}
was not found in the Thomistic corpus by the method used, that Jerome wrote no
homily and no \latin{commentariolus} on the Offertory's psalm at all and
nothing on the Introit's beyond its title and one verse, that Guéranger prints
both the Alleluia and the Offertory without a word of comment, that no witness
read for this guide connects any of the three appointed psalms to the ten
lepers or to the Gospel's thanksgiving, and that the word for leprosy does not
occur anywhere in Theodoret's commentary on the hundred and fifty psalms. No
exhaustive search
of Syriac originals, chant manuscripts, untranslated homiliaries, subscription
databases or current specialist monographs is claimed, and no negative here is
offered as a claim about literature outside the corpora named.

\textbf{Bounds carried by particular claims.} Theodoret's modern English
translation is in copyright, so he is summarised and never quoted throughout,
and the typesetting of this guide carries no Greek, so his terms appear
transliterated where they appear at all. On Galatians he was read in the Greek
column of a Migne scan whose parallel Latin column is destroyed in the
available optical transcription, and the columns given for him there are
approximate and flagged. On the Psalms the reverse problem holds --- that
volume was optically read by a Latin engine, so its Greek columns are destroyed
and its Latin columns and running heads are legible --- and the columns printed
here are the markers embedded in two independent digitisations of the Greek,
four of which were matched against the running heads of the page images in a
parallel pass. Cassiodorus was read in an electronic transmission that states
no underlying edition and carries no PL column marker, so no PL 70 column is
asserted for him. Bellarmine's public-domain English states that it omits
whatever is ``purely philological, or which relate[s] to the discrepancy and
reconciliation of the texts and versions'' and replaces his prefaces with the
Douay headings, so every textual point taken from him --- including his defence
of the Gradual's pronoun, his judgment on the Introit's close, and his note
naming the Mass propers as a witness for the Offertory's word --- is quoted
from the 1611 Latin, read on page images because the optical text of that folio
is unusable. The \work{Glossa ordinaria} was read in a modern electronic
edition; its Rusch facsimile was not opened and its psalter-family apparatus is
reported as that edition's, not as this guide's finding.  Augustine's psalm expositions were read in Latin at a text
site and in the public-domain Oxford English where it exists, and their PL
columns are likewise not asserted. Column numbers for Bede come from
a transcription's embedded markers, not from a facsimile; scan-grade column
evidence exists for Bede, for Ambrose's one clause and for Bruno's commentary,
and for the other Latin witnesses the texts were verified in reliable
electronic transmissions without a column-bearing scan being opened. Cyril's
material for this pericope survives only as catena fragments --- the Syriac of
the relevant sermons is lost --- so no integral sermon on the ten lepers is
cited. Titus of Bostra is fragment-level and unaudited against an independent
edition. Theophylact was verified only in the Latin excerpts of the
\work{Catena aurea}; the columns of PG 123 could not be opened. Rupert of Deutz
is quoted only as Guéranger excerpts him; \work{De divinis officiis} was not
opened. The attribution of the Gregorian supplement is not settled here,
because the scholarship that revised the older attribution is not public
domain. Schuster's cross-reference of the Postcommunion to a Lenten weekday is
reported as his statement and was not independently collated. Schuster (1927)
is public domain in the United States and may remain protected in life-plus-70
jurisdictions; he is quoted in short compass with attribution.

\textbf{Psalm numbering.} Three series are in use and two of them are commonly
both called ``modern.'' \textit{Vulg.} is the Vulgate/Septuagint number, which
the missal cites and which the Douay--Rheims carries. \textit{Eng.} is the
convention of the King James and Revised Standard versions, which leave the
inscription unnumbered and so run a verse or two behind; the page-2 dossiers
print it beside the Vulgate number. \textit{Heb.} is the Masoretic number used
by the Nova Vulgata and the New American Bible, which renumbers the psalm but
keeps the Vulgate's verse division: Ps.~30:15--16 Vulg.\ is Heb.\ 31:15--16;
Ps.~73:1, 19--20, 22--23 is Heb.\ 74:1, 19--20, 22--23; Ps.~89:1 is Heb.\ 90:1;
and Ps.~88:51, cited once as apparatus, is Heb.\ 89:51. The conversions were
taken from the concordance registered with the Douay--Rheims edition used
here, not computed by this guide.

\textbf{Historical bounds.} Dates and places in the page-2 dossiers
distinguish traditional attribution from current historical judgment and are
orienting rather than probative. This guide does not settle the year of the
Passion, the date or place of Luke's composition, the North- or South-Galatian
destination of the Epistle, or the historical authorship of any psalm. Psalm
73's superscription names Asaph while its content laments a sanctuary already
destroyed; the tension is reported and not resolved.

\textbf{Rights.} No public-domain conclusion is claimed for the 1962 typical
edition. This repository's rights record for the facsimile finds that edition's
own status unresolved, and more likely than not restored in the United States
under 17 U.S.C.\ \S104A to the end of 2057 (that record is
\texttt{missale-romanum-1962-facsimile-rights-v1.toml}, under
\texttt{src/sources/inventories/}). What permits the Latin quoted here is
17 U.S.C.\ \S103(b): the copyright in a revised edition extends only to what
that edition itself contributed and ``does not affect or enlarge the scope,
duration, ownership, or subsistence of'' any copyright in the preexisting
material. This formulary is preexisting material. All ten appointed elements
--- and, unusually, the five chant readings that diverge from the Clementine
--- are printed in the Pustet printing of Ratisbon, 1862, which this repository
tracks as public-domain text, and all ten were identified by incipit and order
in the tracked Venice Missal of 1570 besides; the element-by-element check and
its bounds are recorded in \texttt{propers/verified.md}. That is this project's
own reading of the statute and of the evidence, not a clearance anyone granted,
and it is not legal advice. No page image of either digitisation is reproduced
here, and none of the 1962 edition's own new matter is published in this guide;
the edition's rank, pagination and marginal numbers are cited as facts about
the book. The project claims no rights in the received liturgical text. Both
digitisations used to verify it are third-party artifacts whose own rights
status the repository's source records list as unresolved, and neither is
redistributed here. Quotations from sources still in copyright are avoided or
kept to short compass with attribution, as noted at each point of use. The
common reuse notice on the final page states the outbound boundary and does not
relicense any of this material.

\textbf{Review state.} This document has had internal production review only:
source collation, build, and page-by-page visual inspection by the agents that
produced it. It has received no independent liturgical, theological,
specialist, rights or ecclesiastical review, and carries no imprimatur, nihil
obstat or ecclesiastical approval. The exact published PDF snapshots have
distribution authorization under the release record; that authorization does
not relicense the incorporated liturgical text for independent extraction and
does not constitute any of the outstanding reviews.
